\FloatBarrier
\subsection{Development Audit of Pool Expansion}
\label{sec:pool_expansion_audit}

The purpose of the pool-expansion audit is operational rather than
probabilistic. It asks whether the module can turn zero-capacity control cells
into selectable households under the released defaults, how many records it
needs to do so, and whether the selected records satisfy the same rule set as
the rest of the pool. It does not claim that every emitted record is realised
in the census, nor that the expansion closes the zone's full deficit; those are
questions for the zone fit and the validation tables.

\paragraph{Setting}
The audit uses ten SA2 zones from the SEQ case study, the CP feasible pool
(1{,}396{,}307 rows), the 1{,}252{,}479-row census extract, and the released
defaults: \texttt{rep\_max}~$=40$, a per-cell carrier requirement
$\nu_c=\lceil E_{z,k,m}/r_{\max}\rceil$ and a per-zone budget
$K_z=\sum_c\nu_c$ households ($P_z^{\mathrm{add}}=M K_z$ persons), four donors
per zero cell, and a 30-second limit per selection MIP.
The optional support block is \eqref{eq:mip_joint}. A cell is a \emph{zero cell} when
its census target is positive and its attainable capacity in the zone's
candidate pool is zero; cells with positive but insufficient capacity are
partial deficits and are left to the zone fitter, which can replicate a
candidate up to \texttt{rep\_max}.

\paragraph{Coverage}
Of 146 deficient control cells across the ten zones, 123 were zero-capacity and
23 were partial. The module generated donor--shell candidates for 70 of the 73
sourceable zero cells and the joint selection covered all 70. The remaining 53
zero cells were not filled: 40 belong to four control codes that are absent from
the census extract (\texttt{hrsp:101}, \texttt{hrsp:103}, \texttt{qalfp:16},
\texttt{strd:9}), 10 are the \texttt{tend:6} padding code, and three are
\texttt{tenlld:11} cells for which no donor--shell candidate survived the
household and rule checks. The first 50 are cells the current policy treats as absent or padding
rather than as donor shortages; the last three identify a candidate-generation
gap rather than a structural zero. Counting only cells the current policy treats as sourceable,
coverage is 70/73 (95.9\%); treating \texttt{tend:6} as sourceable as well gives
70/83 (84.3\%). Table~\ref{tab:pool_expansion_audit} gives the per-zone result.

\paragraph{Selected records}
The joint selection chose 14 households containing 18 person rows across the
ten zones, one or two households per zone, and never more than three persons.
The board predates both the carrier requirement and the per-SA2 tagged menu
growth; it describes the earlier one-carrier program, so these counts must not
be read as the current program's result. Once every sourceable zero cell has
its carriers, the marginal term does not justify adding more records, and the
selection stops. All ten MIPs returned an optimal solution.
Every selected row passed the eleven consistency-rule checks; the added
households have 14 unique identifiers, no shell household received more than one
transplant, and the output grew from 1{,}396{,}307 to 1{,}396{,}325 rows.
Attributes outside the person-level model were carried by the shell or by a
verbatim donor household, not repaired.

\paragraph{Joint-marginal deviation}
For each zone the weighted deviation
$\sum_{k,m}w_{k,m}d_{k,m}$ of the selected set ranged from 0.239 to 1.622
(median 0.406), and the unweighted mean absolute cell deviation from 0.126 to
0.280 (median 0.147); the largest single-cell deviation in a zone ranged from
0.998 to 1.931. These are deviations of the added set from the zone's target
shares, not measures of the fitted zone population. Because the added set
contains only one to three persons per zone, a single person can move a cell
share by a large amount; the values are reported to make the joint trade-off
visible, not to claim that an 18-person addition reproduces the zone marginals.
The joint model is doing what it is intended to do: it chooses one set that pays
every table's deviation against coverage and against the household count,
rather than repairing each zero cell independently.

\paragraph{What the audit supports}
The audit supports three operational claims. First, the mechanism fills zero
cells that the CP pool cannot supply: all 70 sourceable cells with an accepted
candidate were covered, and the selected records are rule-feasible. Second, it
does so with a very small added set, 14 households and 18 persons, far below
the budget, so feasibility is achieved without forcing the pool to reproduce
the census shortfall. Third, the failures are legible: 50 of the 53 residual
cells are absent or padding codes, and the other three are a specific
donor--shell case to fix. It does not support a claim that every zero cell is
fillable, that the added records are realised in the census, or that the zone
fit is improved by the expansion alone. Those claims require the zone-level
fitting and validation reported in the main text.

\begin{table}[h!]
\centering
\caption{Ten-zone development audit of pool expansion under the released
defaults. \emph{Zero} is the number of positive-target control cells with zero
attainable capacity; \emph{cov.} is the number covered by the selected set;
\emph{HH} and \emph{persons} are the selected added records; \emph{dev.} is the
weighted joint-marginal deviation over the 19 fitting tables
$\sum_{k,m}w_{k,m}d_{k,m}$ (lower is better). All ten support solves
returned an optimal solution; the audit took 1{,}013.5 seconds for the ten
zones. The 53 uncovered cells are 40 cells in four census-absent codes, 10
cells in the \texttt{tend:6} padding code, and three \texttt{tenlld:11} cases
with no valid candidate.}
\label{tab:pool_expansion_audit}
\small
\begin{tabular}{@{}lrrrrr@{}}
\toprule
Zone & Zero & Cov. & HH & Persons & Dev. \\
\midrule
112011236 & 13 & 8 & 2 & 3 & 1.113 \\
112011237 & 13 & 7 & 2 & 2 & 0.239 \\
112011238 & 11 & 6 & 1 & 1 & 0.458 \\
112011239 & 12 & 7 & 1 & 1 & 0.309 \\
112011240 & 13 & 8 & 2 & 3 & 0.355 \\
112011241 & 11 & 6 & 1 & 1 & 0.341 \\
112011242 & 13 & 7 & 1 & 1 & 0.827 \\
112011243 & 12 & 7 & 1 & 1 & 0.772 \\
112021244 & 13 & 8 & 2 & 3 & 1.622 \\
112021245 & 12 & 6 & 1 & 2 & 0.288 \\
\midrule
Total  & 123 & 70 & 14 & 18 & 6.326 \\
Median & 12  & 7  & 1  & 1  & 0.406 \\
\bottomrule
\end{tabular}
\end{table}

